\section{Methodology}

Our goal is to detect when classifier commitment to spurious features accelerates---the crystallization point. We approach this as a signal detection problem: identify a characteristic signature in training dynamics that marks the transition.

\subsection{Overview}

The crystallization zone represents a phase where worst-group accuracy (WGA) decline accelerates. While WGA decreases throughout early ERM training, the rate of decline is not constant. We hypothesize that crystallization manifests as a localized acceleration---a negative second derivative---rather than gradual monotonic decline.

Our detection pipeline consists of three stages: (1) dense WGA tracking during training, (2) smoothing to reduce noise, and (3) second derivative computation with peak detection.

\subsection{Worst-Group Accuracy Tracking}

We compute worst-group accuracy at every epoch during training. For a dataset with group annotations (spurious attribute $\times$ label), WGA is the minimum accuracy across all groups:
\begin{equation}
\text{WGA} = \min_{g \in \mathcal{G}} \frac{1}{|D_g|} \sum_{(x,y) \in D_g} \mathbf{1}[\hat{y}(x) = y]
\end{equation}

Dense checkpointing (every epoch) is essential. Coarser intervals may miss the crystallization peak or shift its apparent timing.

\subsection{Smoothing}

Raw WGA curves exhibit epoch-to-epoch noise due to batch sampling variance. We apply a rolling average with a window of 5 epochs. The choice of 5-epoch window balances noise reduction with temporal resolution. Our sensitivity analysis confirms that windows of 3, 5, or 7 epochs all achieve high detection rates, with 5 epochs providing optimal signal-to-noise ratio.

\subsection{Second Derivative Detection}

We compute the second derivative using central differences on the smoothed WGA. A significant negative peak in $d^2\text{WGA}/dt^2$ indicates accelerating decline---the crystallization signature. We detect peaks using prominence threshold ($>0.005$), search within first 50\% of training, and SNR threshold ($>2.0$).

\subsection{Linear Probe Analysis}

To verify that crystallization represents classifier commitment rather than representation degradation, we analyze feature representations using linear probes on frozen penultimate layer activations.

\subsection{Gradient Ratio Tracking}

To verify the causal role of gradient starvation, we track the gradient ratio between minority and majority group examples. If gradient starvation causes crystallization, we expect the gradient ratio inflection to \emph{precede} the WGA crystallization peak.
